\section{LDP Walk Counting}
\label{sec:k-wk}


We begin with a special case of acyclic pattern counting: $k$-line path counting. 
As a preliminary step, we develop a mechanism for the related task of $k$-line walk counting.
% Before that, we first discuss $k$-line walk counting. 
As noted in Section~\ref{subsec:prelim_pattern_counting}, walk counting can be viewed as a simpler variant of path counting.
Although walk counting does not fall under acyclic graph pattern counting, it forms the foundation for our mechanisms for both path counting and general acyclic graph pattern counting.
{In this section, we first provide a brief overview of the existing one-round RR-based solution for walk counting, which incurs an error of $O(N^k)$.}
% In this section, we first present a strawman solution for walk counting using RR, which incurs a variance of $\Omega(N^{k+1})$. 
We then introduce our multi-round LDP approach, which achieves an error bound of $\tilde{O}(Nd(G)^{k-1})$ with constant probability, while improving communication and computation costs by an order of magnitude.

 

% \subsection{Strawman: Releasing Noisy Graph with RR}
\subsection{{Overview of the RR-Based Method}}
\label{subsec:strawman}


% Similar with other most pattern counting tasks under local-DP (e.g., triangle counting), the key challenge, for answering ...,  lies in capturing correlations among multiple edges.
% One straightforward idea to address is issue is to first construct 
% a noisy graph $\tilde{G}$ is constructed by perturbing each edge... and then counting the ... over the noisy graph...
% More precisely, we utilize 



Similar to many pattern counting tasks under LDP, such as triangle counting, the key challenge in answering $k$-line walk queries lies in capturing correlations across multiple edges. 
{\cite{suppakitpaisarn2025counting} adopts a straightforward one-round approach that first constructs a noisy graph $\tilde{G}$ using RR, as described in Section~\ref{subsec:ldp}.
Given $\tilde{G}$, for each node pair $(v_i, v_j)$, the analyzer $\mathcal{A}$ then computes an unbiased estimator $\hat{a}_{i,j}$ for the binary existence of edge $(v_i, v_j)$ in the original graph $G$.
Since the existence of any walk can be unbiasedly estimated by the product of the estimators of its constituent edges, $\mathcal{A}$ enumerates all $(k+1)$-node sequences and sums the corresponding walk estimators to compute the total count.
This method guarantees an unbiased estimate with an additive error of $O(N^k)$ for $k$-line walk counting, and it extends naturally to general $k$-edge acyclic pattern counting with the same asymptotic utility guarantee.
}




\paragraph{Redundant count removal}
The above solution treats walk orientations distinctly; that is, the sequences $(v_1, \dots, v_{k+1})$ and $(v_{k+1}, \dots, v_1)$ are considered different walks. 
{To eliminate redundant counts due to orientation, ~\cite{suppakitpaisarn2025counting} divides the total count by the number of automorphisms~\cite{godsil2013algebraic} of the queried pattern.}
{Due to space constraints, we defer detailed discussions on redundant count handling across pattern counting queries to Appendix~{A.1} of the full version~\cite{full_version}.}





\paragraph{Communication and computation analysis}
The total communication cost is $O(N^2)$, as each node $v_i$ sends a message of size $O(N)$ to construct the noisy graph. 
For computation, the analyzer $\mathcal{A}$ requires $O(N^{k+1})$ time to enumerate
all possible sequences,
% all possible sequences $w$,
while each node incurs $O(N)$ time to perturb its neighbor list.





\subsection{Our Approach: Multi-Round Aggregation}
\label{subsec:k-wk}



The problem with the strawman solution is that we publish ``too much'' information 
of nodes
in an overly ``fine-grained'' manner. 
Fundamentally, the more information is released with the same privacy level, the more difficult it is to achieve higher utility based on the released data.
Therefore, we would like to collect the data in a more ``local'' and more ``aggregated'' manner. 
Given edge counting (i.e., $1$-line walk counting) as an example, instead of letting each node publish its perturbed neighbor list, we release only the degree information.
The analyzer then aggregates the perturbed degrees to produce the query result.
With this approach, the error can be reduced from typically $\Omega(N)$
to $O(\sqrt{N})$, by simply adding Laplace noises of scale $2/\varepsilon$ to each node's degree count. 




\begin{algorithm}
\small
\caption{Local randomizer $\mathcal{R}^{(\ell)}$ of $\mathcal{M}_{\mathrm{k\text{-}wk}}$}
\label{alg:k-wk_randomizer}
\SetNoFillComment

\tcc{Initialization} 
\ForEach{$i \in \{1, \dots, N\}$}{
    $\mathcal{X}^{(0)}_i \leftarrow 1$
}

\SetKwFunction{r}{$\mathcal{R}^{(\ell)}$}
\SetKwProg{Fn}{Function}{:}{}
\Fn{\r{$G_i$, $\varepsilon$}}{

Received $\mathcal{X}^{(\ell -1)}_i$ for each $i \in \{1, \dots, N\}$

    $ \|\mathcal{X}^{(\ell-1)}\|_\infty := \max_i|\mathcal{X}_i^{(\ell-1)}|$ \\
    % $GS_Q \leftarrow \|\mathcal{X}^{(\ell-1)}\|_\infty$\\
    $\mathcal{X}^{(\ell)}_i \leftarrow 
    \sum_{j \in \mathcal{N}(i)}\mathcal{X}^{(\ell-1)}_j +
    \operatorname{Lap}(\frac{2k}{\varepsilon}\|\mathcal{X}^{(\ell-1)}\|_\infty)$\\
    
    \textbf{Publish} $\mathcal{X}^{(\ell)}_i$
   
}
\end{algorithm}






This idea, {also explored in~\cite{betzer2024publishing}}, can be extended to support arbitrary $k$-line walk counting queries with multiple rounds of communication: 
At the beginning, we initialize $\mathcal{X}^{(0)}_i =1$ for any $v_i\in V$, representing that the number of $0$-line walk ending at each node is $1$.
Then, in round $\ell \in \{1, \dots, k\}$, each node collects the number of $(\ell-1)$-line walks ending at its neighbors.
By aggregating these values, each node can compute the number of $\ell$-line walks that end at itself.
To ensure the reported counts satisfy edge-LDP, each node must add noise to obscure the presence or absence of edges between its neighbors.
Specifically, {unlike~\cite{betzer2024publishing} that limits the global sensitivity using a clipping mechanism,}
our approach adds Laplace noise with a scale proportional to the maximum count from the previous round.
% Laplace noise is added with a scale proportional to the maximum count from the previous round.
\uline{Notably, for each node, the noise scale is determined by the maximum count across all nodes rather than just its own neighborhood, thereby preventing leakage of local structural information.}
Besides, since each edge contributes to the computation in every round, the total privacy budget is divided across the $k$ rounds, allocating $\varepsilon/k$ to each.
Finally, the counts produced in round $k$ are sent to the analyzer $\mathcal{A}$, which aggregates all node reports and returns 
\begin{equation*}
\label{e:k_owk_output}
     \mathcal{M}_{\mathrm{k\text{-}wk}}(G,\varepsilon) =  \sum_{i=1}^N \mathcal{X}^{(k)}_i
\end{equation*}
as the output.
The detailed local randomizer $\mathcal{R}^{(\ell)}$ applied at round~$\ell$ of this $k$-round local-DP mechanism $\mathcal{M}_{\mathrm{k\text{-}wk}}$ is described in Algorithm~\ref{alg:k-wk_randomizer}. 



Figure~\ref{fig:2-wk} illustrates the execution of the $2$-round mechanism $\mathcal{M}_{\mathrm{2\text{-}wk}}$ on graph $G$ with a total privacy budget of $\varepsilon = 4$.
Each node is assigned a per-round privacy budget of $\varepsilon / 2k = 1$.
In both rounds $1$ and $2$, each node first aggregates the outputs from its neighbors in the previous round (shown in blue), and then adds Laplace noise to privatize the result (shown in red).
The final output of the mechanism, $5.7$, is computed by summing the privatized values of all nodes in round $2$.



\begin{figure}[h]
  \centering
  \includegraphics[width=0.58\linewidth]{figures/2-wk.pdf} % file myfig.svg in figs/
\caption{Illustration of the 2-line walk counting mechanism $\mathcal{M}_{\mathrm{2\text{-}wk}}$ applied to graph $G$ with privacy budget $\varepsilon = 4$.
For simplicity, $\sum$ denotes $\sum_{j \in \mathcal{N}(i)}$.}
\label{fig:2-wk}
\end{figure}



For privacy and utility, we establish the following guarantees:
\begin{theorem}
\label{the:k-owk-privacy}
For any graph $G$, walk length $k$, and privacy budget $\varepsilon$, $\mathcal{M}_{\mathrm{k\text{-}wk}}$ satisfies $\varepsilon$-edge-LDP.
\end{theorem}


\begin{theorem}
\label{the:o_k-wk_utility}
For any graph $G$, walk length $k$, and parameters $\varepsilon$ and $\beta$, the $k$-round mechanism $\mathcal{M}_{\mathrm{k\text{-}wk}}$ satisfies
$\mathbb{E}[\mathcal{M}_{\mathrm{k\text{-}wk}}(G,\varepsilon)] = {W}_k$,
and with probability at least $1 - \beta$,
\begin{equation*}
|\mathcal{M}_{\mathrm{k\text{-}wk}}(G,\varepsilon) - {W}_k|
  \leq  k\gamma \sqrt{N} (d(G) + \gamma)^{k-1} = \tilde{O}(\sqrt{N}d(G)^{k-1}),
\end{equation*}
where $\gamma = 2k\sqrt{8\log(2kN/\beta)}/\varepsilon = \tilde{O}(1)$.  
\end{theorem}


{Proofs of all lemmas and theorems in this paper are provided in Appendix~A.3 of the full version~\cite{full_version}.}
By Theorem~\ref{the:o_k-wk_utility}, $\mathcal{M}_\mathrm{k\text{-}wk}$ successfully reduces the error from typically {$O(N^k)$} to $\tilde{O}(\sqrt{N}d(G)^{k-1})$ for general $k$-line walk counting queries.



\paragraph{Communication and computation analysis}
The total communication cost of the $k$-round mechanism $\mathcal{M}_{\mathrm{k\text{-}wk}}$ is $O(N^2)$. 
Specifically, in the first $k-1$ rounds, each node $v_i$ broadcasts its count of $O(1)$ size to all other nodes, incurring a total cost of $O(N^2)$, considering that $k$ is a constant. 
In the final round, each node sends its count to the analyzer, contributing an additional $O(N)$ cost.
For the computational cost, the analyzer $\mathcal{A}$ runs in $O(N)$ time to aggregate the values received from all $N$ nodes in round $k$. 
Each node performs $O(N)$ computation across all $k$ rounds: in each round, it spends $O(N)$ time computing a maximum and another $O(N)$ time aggregating values from neighbors.



\subsection{Optimization in Efficiency}
\label{subsec:opt_wk}


So far, $\mathcal{M}_{\mathrm{k\text{-}wk}}$ has achieved improved computational complexity with state-of-the-art utility. 
Next, we further optimize both communication and computation.


\paragraph{Round reduction:} 



The number of rounds in $\mathcal{M}_{\mathrm{k\text{-}wk}}$ can be reduced from $k$ to $k-1$. 
The key observation is that, for each $k$-line walk, the final edge is shared between the second-to-last node and the last node. 
Therefore, once the second-to-last node obtains the noisy count of $(k-1)$-line walks ending at itself, it can compute the number of $k$-line walks where it serves as the second-to-last node, eliminating the need for round $k$.
More precisely, in round $k-1$, each node $v_i$ computes $\mathcal{X}^{(k-1)}_i$ as in Algorithm~\ref{alg:k-wk_randomizer} and multiplies it by $d(v_i) + \mathrm{Lap}(2k/\varepsilon)$. 
The analyzer $\mathcal{A}$ then aggregates these values to produce the final output. 
We will show that this modification preserves both the privacy and utility guarantees stated in Theorem~\ref{the:k-owk-privacy} and Theorem~\ref{the:o_k-wk_utility}.



\paragraph{Computational cost reduction:} 




Another observation is that, in each round, each node scans through all $N$ results from the previous round to compute a maximum, which incurs significant computational redundancy. 
To address this, we can shift the maximum computation to the analyzer. 
More precisely, 
at the end of each round $\ell \in \{1, \dots, k-2\}$, each node $v_i$ additionally sends its value $\mathcal{X}^{(\ell)}_i$ to the analyzer $\mathcal{A}$, which computes the global round maximum and broadcasts it to all nodes. 
This reduces the computation cost at each node from $O(N)$ to $O(d(G))$.




\paragraph{Communication cost reduction:}

The communication cost between nodes can be reduced from $O(N^2)$ to $O(M)$ under the assumption that nodes connected by an edge have direct communication channels. 
In this setting, at the end of each of the first $k-2$ rounds, each node sends its result only to its neighbors, while the global maximum is computed by the analyzer as discussed above. 
This change does not affect the node-level computation, as each node only requires the results from its neighbors and the global maximum. 
Consequently, in each of the first $k-2$ rounds, each edge facilitates bidirectional communication between its endpoints, leading to a total communication cost of $O(M)$ between nodes.





\begin{algorithm}
\small
\caption{Local randomizer $\mathcal{R}^{(\ell)}$ of optimized $\mathcal{M}_{\mathrm{k\text{-}wk}}$}
\label{alg:k-wk-opt_randomizer}
\SetNoFillComment

\ForEach{$i \in \{1, \dots, N\}$}{
    $\mathcal{X}^{(0)}_i \leftarrow 1$
}

\SetKwFunction{r}{$\mathcal{R}^{(\ell)}$}
\SetKwProg{Fn}{Function}{:}{}
\Fn{\r{$G_i$, $\varepsilon$}}{

% \tcc{Computational cost reduction}
Received $\|\mathcal{X}^{(\ell-1)}\|_\infty$ from $\mathcal{A}$

% \tcc{Communication cost reduction}
   Received $\mathcal{X}^{(\ell-1)}_j$ for each $j \in \mathcal{N}(i)$

    
    $\mathcal{X}^{(\ell)}_i \leftarrow 
    \sum_{j \in \mathcal{N}(i)}\mathcal{X}^{(\ell-1)}_j +
    \operatorname{Lap}(\frac{2k}{\varepsilon}\|\mathcal{X}^{(\ell-1)}\|_\infty)$\\

  \If{$\ell = k-1$}{
        % \tcc{Round reduction}
         ${\mathcal{X}}^{(\ell)}_i 
        \leftarrow 
        {\mathcal{X}}^{(\ell)}_i \cdot 
        (d(v_i) + \operatorname{Lap}(\frac{2k}{\varepsilon}))$\\
    }

    
    \textbf{Send out} $\mathcal{X}^{(\ell)}_i$ to $\mathcal{N}(i)$ and $\mathcal{A}$
   
}
\end{algorithm}



The revised local randomizer is shown in Algorithm~\ref{alg:k-wk-opt_randomizer}.
These refinements collectively reduce the total number of rounds from $k$ to $k-1$, lower the overall communication cost from $O(N^2)$ to $O(M+N)$, and achieve an
improvement from $O(N)$ to $O(d(G))$ in the per-node computational cost.
Below, we show that these refinements will not affect the privacy and utility guarantee:
\begin{theorem}
For any graph $G$, walk length $k$, and privacy budget $\varepsilon$, the optimized mechanism $\mathcal{M}_{\mathrm{k\text{-}wk}}$ satisfies $\varepsilon$-edge-LDP.
\end{theorem}


\begin{theorem}
For any graph $G$, walk length $k$, and parameters $\varepsilon$ and $\beta$, the optimized $(k-1)$-round mechanism $\mathcal{M}_{\mathrm{k\text{-}wk}}$ satisfies
$
     \mathbb{E}[\mathcal{M}_{\mathrm{k\text{-}wk}}(G,\varepsilon)] = {W}_k,
$
and with probability at least $1 - \beta$,
\begin{equation*}
|\mathcal{M}_{\mathrm{k\text{-}wk}}(G,\varepsilon) - {W}_k|
  \leq  k\gamma \sqrt{N} (d(G) + \gamma)^{k-1} = \tilde{O}(\sqrt{N}d(G)^{k-1}),
\end{equation*}
where $\gamma = 2k\sqrt{8\log(2kN/\beta)}/\varepsilon = \tilde{O}(1)$.  
\end{theorem}





% move to former part.......
% \subsection{Unoriented Walk Counting}



% In the earlier part of this section, we focused on counting oriented walks, where two $k$-walks—$w_1$ from $v_1$ to $v_{k+1}$ and $w_2$ from $v_{k+1}$ to $v_1$—are treated as distinct. 
% This reflects common scenarios where a trip from location $a$ to $b$ is considered different from one in the reverse direction. 
% Let $W_k$ denote the number of unoriented $k$-walks of the input graph. 
% Since walks may contain repeated nodes or edges, a walk can be symmetric, that is, its node sequence remains the same under reversal.
% For example, $(v_1, v_2, v_1)$ is a symmetric $3$-line walk.
% As a result, $W_k$ cannot be derived directly from having the oriented walk count $O_k$, because symmetric walks are counted only once in $O_k$.
% Nevertheless, both the strawman and our proposed mechanisms can be extended to handle the unoriented walk counting problem under local-DP, where $w_1$ and $w_2$ are treated as equivalent.


% % \textbf{Our extension.}
% We construct our local-DP scheme $\mathcal{M}_{\mathrm{k\text{-}wk}}$ for counting unoriented $k$-line walks, building upon $\mathcal{M}_{\mathrm{k\text{-}owk}}$ as follows.
% The key idea is that although $W_k$ cannot be directly derived from $O_k$, it satisfies the relation $W_k = (O_k + U_k)/2$, where $U_k$ denotes the number of symmetric walks of length $k$. 
% Since $\mathcal{M}_{\mathrm{k\text{-}owk}}$ produces an unbiased estimator of $O_k$, the task of estimating $W_k$ reduces to estimating $U_k$.
% Intuitively, a symmetric $k$-line walk is fully determined by its first $\lceil k/2 \rceil$ nodes, with the remaining nodes mirrored accordingly. 
% Hence, $U_k = O_{\lceil k/2 \rceil}$, further reducing the task to estimating $O_{\lceil k/2 \rceil}$.
% Recall that $\mathcal{M}_{\mathrm{k\text{-}owk}}$ computes, in each round $\ell$ from $1$ to $k$, the number of oriented $\ell$-line walks ending at each node.
% Thus, we construct $\mathcal{M}_{\mathrm{k\text{-}wk}}$ by introducing an additional step at round $\lceil k/2 \rceil$, in which the analyzer not only computes the maximum of the received node counts, but also aggregates them to compute an unbiased estimator $\hat{U}_k$ of $O_{\lceil k/2 \rceil}$.
% In the final round, the analyzer outputs $(\hat{O}_k + \hat{U}_k)/2$, yielding an unbiased estimator of $W_k$.

% In general, this modification adds an additional $O(N)$ computational cost while leaving the communication cost unchanged, neither of which asymptotically affects the overall cost.
% Since the additional step involves only post-processing of privatized results, $\mathcal{M}_{\mathrm{k\text{-}wk}}$ also satisfies $\varepsilon$-edge-DP.
% In terms of utility, the overall error is bounded by the sum of the error bounds for $\hat{O}_k$ and $\hat{U}_k$. 
% The error in $\hat{O}_k$ asymptotically dominates, so the overall utility bound remains $\tilde{O}(\sqrt{N}D^{k-1})$.
% Compared to the centralized-DP Laplace mechanism, which incurs error $O(D^{k-1})$ on the same query, our method still achieves the optimal amplification factor of $O(\sqrt{N})$.



% \green{here}



% example? 













